\section*{Chapter \ref{ch:iv3:operating-systems-concurrency-and-networks} --- Operating Systems, Concurrency and Networks}
\markboth{\protect\omTocMark{}{Answers: Operating Systems, Concurrency and Networks}}{\protect\omTocMark{}{Answers: Operating Systems, Concurrency and Networks}}

\begin{solution}{iq:iv3:operating-systems-concurrency-and-networks:1}
It prints \texttt{after fork 0} and \texttt{after thread 1}. The child has its own copy of \texttt{count}, so its increment stays in the child, while the thread changes the one \texttt{count} in the process.

Threads in a process share the address space, so they share global data and the heap. They also share open files, and signal handlers. Each thread has its own stack, its own registers and program counter, and its own thread-local storage. A thread can still reach another's stack through a pointer, as the memory is all one space.

After \texttt{fork}, the child starts with a copy of the parent's memory, which Linux makes lazily, and a copy of its file descriptor table. No memory stays shared unless it was mapped as shared. The open file descriptions are the exception: parent and child share file offsets through them.

\iqwrong{``A process and a thread differ only in cost.'' The cost follows from the real difference, which is whether writes are visible to the other side.}
\iqfollow{Which is safer when one side crashes?}{Processes: a wild write in one cannot corrupt the other's memory, and a crash kills only that process. A segmentation fault in one thread ends all of them.}
\end{solution}

\begin{solution}{iq:iv3:operating-systems-concurrency-and-networks:2}
It prints \texttt{4000 1000 0}. \texttt{shared} is one atomic that all four threads increment. \texttt{mine} is one variable per thread, so each thread counts to 1000, and the \texttt{mine} of the main thread, which did no counting, is still 0.

\texttt{shared} has one address, in the program's data segment. Each thread gets its own block of thread-local storage when it starts, and \texttt{mine} sits in it at a fixed offset. On x86-64 Linux, a register (\texttt{fs}) points at the current thread's block, so reaching \texttt{mine} takes no lock and needs no atomic.

That suits a statistic updated often and read rarely. Each thread increments a \texttt{thread\_local} counter of its own, which no other core writes, and a reader adds them up from time to time. That avoids a cache line bouncing between cores for each increment, at the price of a sum that is not a snapshot, and the counters must be registered so that a reader can find them.

\iqwrong{``\texttt{thread\_local} is the same as a local variable.'' A local one is new on each call, whereas this lasts as long as its thread does and keeps its value between calls.}
\end{solution}

\begin{solution}{iq:iv3:operating-systems-concurrency-and-networks:3}
It prints \texttt{state Z}, \texttt{code 7} and \texttt{gone True}. A zombie is a child that has exited and not yet been waited for. It holds no memory or open files, only its entry in the process table, which keeps the exit status for the parent.

The way to clear one is for the parent to call \texttt{wait} or \texttt{waitpid}, as the listing does, and that gives the parent the code 7. You cannot kill a zombie, as it is already dead, so \texttt{kill -9} does nothing.

If the parent dies first, its children are handed to \texttt{init} or the nearest subreaper, and that process waits on them, so the zombies go. The harm comes from a long-lived parent that never waits: each zombie holds a process ID, and the system has a limited number of these.

A parent that does not care about the status can set \texttt{SIGCHLD} to be ignored, and then no zombie is left.

\iqwrong{``A zombie is a process that hangs and uses CPU.'' It uses no CPU, as it has already exited.}
\end{solution}

\begin{solution}{iq:iv3:operating-systems-concurrency-and-networks:4}
It prints \texttt{b\textquotesingle{}abc\textquotesingle{} b\textquotesingle{}de\textquotesingle{}} and \texttt{b\textquotesingle{}abcde\textquotesingle{}}. The number of \texttt{recv} calls for TCP is not fixed: two or three here, counting the last one that returns nothing for the close, and that depends on timing and on the buffers. UDP keeps message boundaries, and TCP is a byte stream with none.

TCP guarantees that the bytes arrive complete, in order and without duplicates, or that the connection reports failure. It adds retransmission, flow control and congestion control, and it needs a connection first. UDP sends each datagram on its own, which may be lost, duplicated or reordered, and it tells the sender nothing. A datagram that does arrive arrives whole.

The trading angle is the cost of order. A lost TCP segment holds back every later byte until it is resent, and a market data consumer would rather have the newer prices now. So feeds usually send UDP with sequence numbers, and handle gaps themselves, while order entry usually runs over TCP, as each message must arrive and in order.

\iqwrong{``UDP is unreliable, so it loses data all the time.'' On a healthy network it rarely does. The point is that nothing repairs it when it does.}
\end{solution}

\begin{solution}{iq:iv3:operating-systems-concurrency-and-networks:5}
It prints \texttt{little}, \texttt{04030201 01020304} and \texttt{0x4030201}. x86-64 stores the low byte of a number first, so 0x01020304 is laid out as 04 03 02 01. Network byte order is big-endian, the high byte first, which is what the format \texttt{!I} packs and what \texttt{htonl} converts to.

So \texttt{htonl} reverses the bytes on this machine. On a big-endian machine, it would do nothing.

A feed that sends its prices in little-endian order is read directly with the native format on x86-64, with no reversal. On a big-endian machine, the same price needs an explicit little-endian format. For portable code, write the order you mean: \texttt{<I} says little-endian on any machine, while \texttt{=I} means whatever the machine uses.

\iqwrong{``Network byte order is what the network card uses, so the kernel handles it.'' The kernel converts its own headers. The fields in your payload are yours to handle.}
\iqfollow{What does a wrong order look like?}{A small value turns huge and a huge one turns small, by a factor of up to $2^{24}$ (exactly that for a value that fits in one byte), which should prompt a sanity check on any price.}
\end{solution}

\begin{solution}{iq:iv3:operating-systems-concurrency-and-networks:6}
It prints \texttt{True}, \texttt{False} and \texttt{b\textquotesingle{}a\textquotesingle{} b\textquotesingle{}b\textquotesingle{}}. Both connections have the same local address and port, and different peer ports. A TCP connection is identified by four things: local address, local port, remote address and remote port. With the protocol, that makes the five-tuple.

So one server port can carry many connections. The kernel looks up each arriving segment by that tuple, and hands it to the right socket. \texttt{accept} returns a new socket for each connection, and the listening socket goes on listening.

The real limits are elsewhere. A client talking to a single server address and port has only its range of ephemeral source ports, some tens of thousands here (\texttt{ip\_local\_port\_range}), and a server's limit comes from memory and file descriptors.

\iqwrong{``Each client needs its own server port.'' That confuses the listening port with the connection. No new port is opened on \texttt{accept}.}
\iqfollow{How is a UDP socket different?}{It has no connection, so one socket receives from all senders, and the source address comes with each datagram.}
\end{solution}

\begin{solution}{iq:iv3:operating-systems-concurrency-and-networks:7}
It prints \texttt{1110} and \texttt{10}. Three slots can be taken and the fourth try fails. After the other thread's release, one slot is free again. A semaphore is a counter with no owner, so any thread may release it, whereas a mutex has an owner, and only that thread may unlock it.

So a semaphore suits counting: at most $N$ requests in flight, or items available in a queue, with one thread acquiring and another releasing. It also signals between threads, which a mutex cannot do. A mutex protects a critical section, and its owner rule lets a system offer priority inheritance, or a debug check.

Unlocking a \texttt{std::mutex} from a thread that does not own it is undefined behaviour. A binary semaphore looks like a mutex, and it gives no such protection, as nothing ties the release to the acquirer.

\iqwrong{``A binary semaphore and a mutex are the same.'' The two have the same states, but only a mutex has an owner.}
\iqfollow{Can \texttt{try\_acquire} fail when a slot is free?}{The standard allows a spurious failure, so a real program should not treat failure as proof that none is free. It did not happen in our runs, which had no contention.}
\end{solution}

\begin{solution}{iq:iv3:operating-systems-concurrency-and-networks:8}
It prints \texttt{caught [15]}, \texttt{refused} and \texttt{-9}. \texttt{SIGTERM} is a request, which a process may handle. \texttt{SIGKILL} cannot be caught, blocked or ignored, as the operator needs one way that always works. The status $-9$ shows the child was killed by that signal.

The kernel simply ends a process on \texttt{SIGKILL}. No handler runs, no destructor runs and no buffered data is flushed.

So a trading process should receive \texttt{SIGTERM} first, and its handler should set a flag and return. The main loop then stops sending new orders, cancels what is open, flushes its log, and exits. Only if it fails to finish in time does the operator send \texttt{SIGKILL}. And since the process can die anyway, in a crash for instance, an exchange-side cancel-on-disconnect, where a venue offers one, is the protection that does not depend on the process.

A handler should do very little. Most library calls are not safe inside one, because the signal may arrive while the same library is holding a lock.

\iqwrong{``Put the cancel logic in the \texttt{SIGKILL} handler.'' There is no such handler. It has to be on \texttt{SIGTERM}, and a crash must be covered elsewhere.}
\end{solution}

\begin{solution}{iq:iv3:operating-systems-concurrency-and-networks:9}
The child takes about 0 faults on reading and about 2000 on writing. In the run here, the two counts were exactly 0 and 2000. The child prints 2 and the parent 1, at the same address, as each process has its own page tables that map one virtual address to different physical pages.

\texttt{fork} copies the page tables, and marks the writable pages read-only in both. So a read hits a page that is already mapped, and no fault occurs. The first write to each page faults. The kernel then copies that page, gives the writer its own copy, and makes it writable. That is copy-on-write. So \texttt{fork} itself is cheap in data, and the cost comes later, a page at a time.

For a large process that forks to save a snapshot, this is a good trick: the child sees a frozen view and writes it out, while the parent keeps running. But each page the parent changes during the save now costs a fault, and a copy of 4 KiB. If the parent rewrites most of its memory, that doubles memory use at worst, and it adds latency to the writes on the hot path, spread unpredictably in time. \texttt{fork} itself must also copy the page tables, which grows with the mapped size, and for that time the process is stopped.

The listing turns huge pages off for this range, so that the parent's first touch does not depend on the system's setting. On a current kernel the child's counts do not change: a write to a shared huge page splits the mapping and copies only 4 KiB, and a run here with huge pages on again gave 0 and 2000. Older kernels tried to copy the whole huge page.

\iqwrong{``The child gets a copy of all the memory at \texttt{fork}.'' It gets the mapping. Physical copies appear on the first write.}
\end{solution}

\begin{solution}{iq:iv3:operating-systems-concurrency-and-networks:10}
The indices are 254, 72, 418 and 359, with offset 2748 (0xabc): bits 47--39, 38--30, 29--21, 20--12, and 11--0. Mapping 1 GiB takes 262\,144 entries, which is 512 tables of 4 KiB, or 2 MiB, plus 3 upper tables: 2\,109\,440 bytes in all. With 2 MiB pages, the last level goes, and it takes 3 tables, 12 KiB.
\omcode{code/interviews-3/05-operating-systems-concurrency-and-networks/python/iv_page_walk.py}{1}{17}{Splitting an address, and counting tables.\protect\hfil\protect\break}\label{lst:iv3:os:pagewalk}

Each table has 512 entries, which is $2^9$, so each level takes 9 bits. A page of 4096 bytes takes 12 bits as the offset. $4 \times 9 + 12 = 48$. On Linux, user space uses the lower half, under 47 bits, which is about 128 TB.

The loop works up the levels. $2^{30} / 2^{12} = 262\,144$ pages need that many entries of the last level, which fill 512 tables; these need 512 entries one level up, in 1 table; and so on up to the top. So the overhead is about 0.2\% of what is mapped. The test checks the split by rebuilding 1000 random addresses from their parts, and the table count against a slow walk over every page that collects the distinct tables. One page alone still needs one table at each level.

The page tables are not what makes huge pages attractive. A translation cache entry then covers 2 MiB and not 4 KiB, so far fewer misses occur, and each miss has one level less to walk.

\iqwrong{``A process's page tables take 8 bytes for each byte.'' It is 8 bytes for each page, which is why the overhead is small.}
\iqfollow{And mapped but untouched memory?}{Linux fills in tables only for pages touched, so an idle mapping costs almost nothing.}
\end{solution}

\begin{solution}{iq:iv3:operating-systems-concurrency-and-networks:11}
It prints \texttt{heap mmap}, \texttt{unmapped} and \texttt{heap}. Small blocks come from the heap, carved up by \texttt{malloc} itself. A block above a threshold, initially 128 KiB, gets its own mapping from the kernel, which \texttt{free} gives back at once, so the second line finds it unmapped.

The last line surprises people. glibc's threshold moves. When a block above the threshold and below its upper limit is freed, the threshold rises to that block's size. So after freeing 1 MiB, the next 1 MiB request comes from the heap. The reasoning is that a program which frees and reallocates big blocks should not pay for a system call every time.

For latency, that matters in three ways. An \texttt{mmap} and a \texttt{munmap} are system calls. A fresh mapping is empty, so each new page faults on first touch. And because the threshold moves, the same request can cost differently depending on history, so a benchmark may not tell you what the first run does.

The cure on a hot path is not to allocate there. Allocate up front, touch the pages, and reuse the memory. To remove the moving threshold, set \texttt{M\_MMAP\_THRESHOLD} with \texttt{mallopt}, which turns off the dynamic behaviour.

\iqwrong{``\texttt{free} returns memory to the operating system.'' Only for a mapped block, or the top of the heap past a trim threshold. Most frees just return it to \texttt{malloc}'s lists.}
\end{solution}

\begin{solution}{iq:iv3:operating-systems-concurrency-and-networks:12}
It prints about \texttt{1024 0} and then \texttt{0 100}. The call gives address space only, with no physical memory behind it. A page gets physical memory the first time it is touched, by a page fault. So the virtual size jumps by 1024 at once, and the resident size grows only for the 100 MiB touched.

The first line is a promise, and the second is reality. The kernel uses this for lazy allocation, for fork's copy-on-write, and for programs that map far more than they use. In its default mode it refuses only obvious overcommits, so a 1 GiB request here succeeds, and a request for far more than the machine could provide would fail. The exact rule depends on a system setting.

If every process touches all it asked for and memory runs out, the failure does not come from \texttt{mmap}, which returned long ago. It comes at a later touch, where the kernel has no pages to give, and it picks a victim with the out-of-memory killer. So a process can die from an ordinary write, on a line with no allocation.

For a trading process, that suggests allocating and touching memory at start-up, and locking it with \texttt{mlock} if a page fault in the hot path cannot be accepted. A failure at start then shows up before the market opens.

\iqwrong{``Checking \texttt{malloc}'s return value protects you from running out of memory.'' It helps with address-space limits. On an overcommitting system, a non-null pointer does not promise any pages.}
\end{solution}

\begin{solution}{iq:iv3:operating-systems-concurrency-and-networks:13}
1000 prints 1000, using about 1 MiB. 100\,000 needs about 100 MiB under an 8 MiB limit, and the process dies with \texttt{SIGSEGV}. Under the 256 MiB limit, it prints 100000. You cannot catch this with \texttt{try}: stack overflow is a memory fault, not a C++ exception. A new thread's stack is 8 MiB under this limit, which our test reads back.

The stack grows downward page by page, as it is touched. When it reaches the limit set by \texttt{ulimit -s}, the kernel refuses to extend it, and the access faults. A signal handler for this would itself need a stack, so it has to run on an alternate one (\texttt{sigaltstack}), and even then it can do little but report.

A new thread does not grow in this way. Its whole stack is mapped when the thread is created, with a size set by \texttt{RLIMIT\_STACK} at program start. If that limit is unlimited, it is 2 MiB on x86-64 instead, so changing \texttt{ulimit} can change every thread's stack size in surprising ways. Set it explicitly with \texttt{pthread\_attr\_setstacksize} when it matters.

That mapping is virtual memory, so 1000 idle threads of 8 MiB take 8 GB of address space, and little real memory. Guard pages lie below each stack, which turns most overflows into a fault, unless a large frame jumps right over the guard.

\iqwrong{``Recursion depth is limited by the language.'' C++ has no such limit. The stack and its setting impose one, and that differs by thread and by machine.}
\end{solution}

\begin{solution}{iq:iv3:operating-systems-concurrency-and-networks:14}
The smallest possible value is 1000: on one core, the child must run between the parent's write and the parent's read returning, so the parent must be taken off the core at least once in each round. It printed 1000 here. A context switch saves one thread's registers, and loads another's, along with its address space if it is in a different process.

The listing adds voluntary switches, when the parent blocks in \texttt{read}, and involuntary ones, when a woken child preempts it, as one can replace the other and only the total has a clean lower bound.

The direct cost is entering the kernel, running the scheduler, and changing registers. Between processes, it also changes the page-table root, which can discard translations. The larger cost usually comes afterwards. The other task has pushed the first one's data out of the caches and its predictions out of the branch predictor, so it runs slowly for a while on return. That cost depends on the work, and does not show up in a plain ping-pong. Published figures vary greatly by machine, so measure yours, and we do not quote one here.

On two cores, with each side spinning on a shared variable and not blocking, neither process leaves its core, so no switch occurs at all. The message crosses by a cache line transfer. The price is two cores at full use, even when idle. That is the usual trade in the hot path of a trading system.

\iqwrong{``A system call and a context switch are the same.'' A system call enters the kernel and usually returns to the same thread, with no switch of task.}
\end{solution}

\begin{solution}{iq:iv3:operating-systems-concurrency-and-networks:15}
Round robin with quantum 2 finishes A at 13, B at 9 and C at 5, for an average of 9. First-come-first-served finishes at 8, 12 and 13, for 11. Shortest-job-first finishes C at 1, B at 5 and A at 13, for 6.33. As the quantum grows, round robin becomes first-come-first-served, which the test confirms.
\omcode{code/interviews-3/05-operating-systems-concurrency-and-networks/python/iv_round_robin.py}{1}{19}{Round robin.\protect\hfil\protect\break}\label{lst:iv3:os:rr}

The trace for quantum 2 goes A to 2, B to 4, C to 5, A to 7, B to 9, and A to 11 and 13. The last completion is always 13, the sum of the work, under any order that never idles. The test checks that, and compares the function with a second tick-by-tick model on 300 random cases.

Shortest-job-first gives the lowest average of the three, and in fact the lowest of any order, since each task's length is counted once for every task behind it. It has two problems. You do not know the lengths in advance, and a steady stream of short jobs starves a long one forever.

Round robin needs no knowledge of length, and it bounds how long anyone waits before running. A short job finishes early, as C does here. A very small quantum gets close to sharing the core equally, but in a real system each switch costs time, which the model ignores, and at some point the quantum is mostly overhead. A real scheduler such as Linux's is more complicated, and tracks how much CPU each task has received.

\iqwrong{``Round robin is fairest, so it is best for average completion.'' On equal-length jobs, it is worse than running them in turn, as all of them finish near the end.}
\end{solution}

\begin{solution}{iq:iv3:operating-systems-concurrency-and-networks:16}
\Cref{lst:iv3:os:waitfor} follows each thread to what it waits on, and then to the holder of that lock, until it repeats itself or stops. Mutual exclusion, hold and wait, no preemption and circular wait must all hold. In practice, you usually remove circular wait, by taking locks in one global order.
\enlargethispage{-2\baselineskip}%
\omcode{code/interviews-3/05-operating-systems-concurrency-and-networks/python/iv_waitfor.py}{1}{17}{Following the waits.\protect\hfil\protect\break}\label{lst:iv3:os:waitfor}

No general search is needed since each thread blocks on at most one lock, and each lock has one holder. So every thread has at most one edge out, and from any start there is just one path to follow. It either stops at a running thread or enters a cycle. \texttt{done} makes the total work linear, as no thread is walked twice.

The other conditions can be removed too. Hold and wait goes if you take all locks at once, as \texttt{std::scoped\_lock} does. No preemption goes if you use a try-lock and back off when it fails, at the risk of livelock.

The test compares the function with a reachability search on 2000 random states, of which over 100 hold a deadlock, and it checks that each returned cycle really is one.

With reader-writer locks or a semaphore, one lock may have many holders. The graph then has several edges per node, and needs a real cycle search, and for a semaphore a cycle no longer proves a deadlock, as another holder may release.

\iqwrong{``A thread blocked for a long time is deadlocked.'' Blocked behind a slow holder is not a deadlock. There must be a cycle, and everyone in it must be waiting.}
\end{solution}

\begin{solution}{iq:iv3:operating-systems-concurrency-and-networks:17}
It prints \texttt{b\textquotesingle{}Saaa\textquotesingle{} b\textquotesingle{}Saaa\textquotesingle{} b\textquotesingle{}SPaa\textquotesingle{}} and \texttt{b\textquotesingle{}SaTa\textquotesingle{} b\textquotesingle{}SaTa\textquotesingle{} b\textquotesingle{}SPaa\textquotesingle{}}. The manual promises that a shared mapping's writes reach the file and other mappings, and that a private mapping's writes stay private. Whether the private mapping sees later changes to the file is unspecified, and that is where the \texttt{S} comes from.

The private mapping got the \texttt{S}, since its page was still the file's own page at that moment, in the page cache. Its first write copied the page. After that it was on its own, and so the later \texttt{T} does not appear. A read through \texttt{pread} sees the shared write at once, as both go through one page cache.

A mapping suits random access over a large file, such as a tick store, since the kernel pages data in on demand, and evicts it under pressure, and it avoids a copy into a user buffer. It also lets several processes share one read-only copy in memory. \texttt{read} is often better for sequential scans. It also reports an I/O error as a return value, whereas a mapped read fails with a signal. And in the hot path, a mapped read can block on a page fault with no system call in sight.

A shared write reaches the disk when the kernel writes back the page, some time later. To force it, call \texttt{msync}, or \texttt{fsync} on the file.

\iqwrong{``\texttt{mmap} is always faster than \texttt{read}.'' Faults have their own cost. For a sequential scan, \texttt{read} with readahead often wins.}
\end{solution}

\begin{solution}{iq:iv3:operating-systems-concurrency-and-networks:18}
\texttt{write} succeeding means the bytes reached the page cache, so the disk may hold the old file, the new one, or any mix. \Cref{lst:iv3:os:atomicsave} writes a new file, flushes it, renames it over the old name, and flushes the directory.
\omcode{code/interviews-3/05-operating-systems-concurrency-and-networks/python/iv_atomic_save.py}{1}{16}{An atomic save.\protect\hfil\protect\break}\label{lst:iv3:os:atomicsave}

Each call has a job. \texttt{f.flush} passes Python's buffer to the kernel. \texttt{fsync} on the file forces the data and its metadata to the device, and it must come before the rename: otherwise a crash might leave the new name on an empty file. \texttt{replace}, which is \texttt{rename(2)}, moves the name in one step, so another process never finds it missing. \texttt{fsync} on the directory makes that rename durable, as fsync on a file does not necessarily save the directory entry.

The test records the order of the calls, and simulates a crash before the rename, which leaves the old contents whole. No test here cuts the power, so the durability itself rests on the manual pages, and on the device.

What remains unprotected: a drive that acknowledges a flush it has not done, which is the hardware's fault. A crash after the rename and before the directory flush can give either version, each complete, and the caller must be able to accept that. Updates between two saves are lost, as this makes each save atomic, not each change. That is why a system that must lose nothing appends to a log, with an fsync per commit, and a save is only an optimisation on top. Each fsync is slow, which is why many systems flush in groups.

\iqwrong{``Rename is atomic, so the fsync on the file is not needed.'' Rename orders the name change, not the data. The name can survive a crash with data that never made it.}
\end{solution}

\begin{solution}{iq:iv3:operating-systems-concurrency-and-networks:19}
Here it prints \texttt{1e-09} and \texttt{0}. Use \texttt{CLOCK\_MONOTONIC} to time a request, as it does not jump, and \texttt{CLOCK\_REALTIME} for the time of day, as only that one relates to a calendar and to other machines. A clock read needs no system call, since the vDSO maps the code and data for it into your process.

\texttt{CLOCK\_REALTIME} can jump when the time is set, so a duration computed from it can come out wrong or even negative. \texttt{CLOCK\_MONOTONIC} does not jump, but its rate is adjusted, slightly, by NTP-like services. \texttt{CLOCK\_MONOTONIC\_RAW} is free of those adjustments, and it is meant for comparison with the hardware. Monotonic time has an arbitrary start, so it says nothing to another machine.

The vDSO exists since such calls are frequent enough to dominate. Reading the clock then costs about a function call, with no kernel entry. We checked that the mapping is there, in \texttt{/proc/self/maps}, and not how each clock source behaves, so check what your machine does.

Resolution is the smallest step that the clock can report, 1 ns here. Accuracy is how close a reading is to true time, and a server may be off by far more than that, depending on how it is synchronised. Two timestamps from different machines can thus be out of order, however fine their resolution.

\iqwrong{``Nanosecond resolution means nanosecond-accurate timestamps.'' The last digits of the wall clock may mean nothing.}
\end{solution}

\begin{solution}{iq:iv3:operating-systems-concurrency-and-networks:20}
It prints \texttt{connected and sent}, then \texttt{b\textquotesingle{}early\textquotesingle{}}. The kernel, not your process, completes the handshake. \texttt{accept} just takes an already established connection off a queue, and data that arrived early waits in that connection's receive buffer.

The handshake goes: the client sends SYN, the server replies with SYN and ACK, and the client sends ACK. \texttt{connect} returns after the second of these. \texttt{listen} tells the kernel to do this for the port. The backlog sets the size of the queue of completed connections waiting for \texttt{accept}, capped by a system setting (\texttt{somaxconn}).

If the queue is full, the manual allows two outcomes: the client gets an error, or its request is ignored so that a retransmission succeeds later. TCP on Linux takes the second by default, so a client may see \texttt{connect} take far longer. A tight server that accepts slowly during a burst of reconnections can stall clients in this way. The behaviour differs by system, and depends on the settings.

So a connection can be established and even hold data, in a server process that has never heard of it. If that process then exits, the client gets a reset. That matters for health checks: a successful TCP connect shows that the kernel is listening, and it does not show that your application is alive.

\iqwrong{``\texttt{accept} performs the handshake.'' It happens beforehand, in the kernel, and the process may be busy or stuck throughout.}
\end{solution}

\begin{solution}{iq:iv3:operating-systems-concurrency-and-networks:21}
It prints \texttt{(True, \textquotesingle{}EADDRINUSE\textquotesingle{})} twice and then \texttt{(True, \textquotesingle{}bound\textquotesingle{})}. The side that closes first enters \texttt{TIME\_WAIT}, here the server's. The second line fails since Linux needs \texttt{SO\_REUSEADDR} on both the old and the new socket: the socket left in \texttt{TIME\_WAIT} was set up without it.

That is what this restart problem usually is. The old process is gone, its connection is in \texttt{TIME\_WAIT}, and the port looks busy for a while. So set \texttt{SO\_REUSEADDR} on the listening socket before \texttt{bind}, from the first version on. It does not allow two active listeners on one address.

\texttt{TIME\_WAIT} has two jobs, and lasts for twice the maximum segment lifetime. The final ACK of the close may be lost, in which case the peer resends its FIN, and this side must still be there to answer. And old segments of this connection may still be in the network, so the same four-tuple must not be reused until they have died out.

So this state is correct, and it costs little: a small kernel object per closed connection. It can become a real problem for a client that opens very many short connections to one server, which can run out of source ports. The right remedy is usually to reuse connections. Setting the state's duration lower or killing it early removes a protection, so be sure you have understood it first.

\iqwrong{``Closing the socket frees the port at once.'' Only if this side was not the active closer, or if it has passed through \texttt{TIME\_WAIT}.}
\iqfollow{Why does the client side not have this problem?}{Usually it does, but its port changes on each connection, so no one notices. Here the server closed first, so it pays.}
\end{solution}

\begin{solution}{iq:iv3:operating-systems-concurrency-and-networks:22}
The product is $10^{10} \times 10^{-3} / 8 = 1.25$ MB. With a 65\,535-byte window, the rate is $65\,535 \times 8 / 10^{-3} = 524$ Mb/s, about 5\% of the link. Window scaling multiplies the 16-bit field by a power of 2, agreed during the handshake, so the window can reach 1 GiB.
\omcode{code/interviews-3/05-operating-systems-concurrency-and-networks/python/iv_window.py}{1}{15}{Window, round trip and rate.\protect\hfil\protect\break}\label{lst:iv3:os:window}

A sender may have at most a window of unacknowledged data in flight, and acknowledgements take a round trip. So one window goes per round trip, whatever the link's speed. To fill the link, the window must reach the bandwidth-delay product. The test compares the formula with a round-by-round count on 200 random cases, and it checks the two ends: a small window limits the rate, and a large window leaves only the link.

Double the round-trip time with a fixed window, and the rate halves. The extreme is an intercontinental link, where a small default window gives little throughput on a very fast path. That's why buffer settings matter more with distance.

The two controls answer different questions. Flow control protects the receiver: it advertises how much buffer it has, and that is the window in the header. Congestion control protects the network: the sender keeps its own congestion window, grows it, and cuts it on signs of loss or delay. The sender uses the smaller of the two.

The answer ignores loss. A lost segment makes a typical sender cut its rate, so on a long path even a small loss rate can reduce throughput far below this limit.

\iqwrong{``A faster link fixes a slow transfer.'' If the window limits the transfer, a faster link changes nothing.}
\end{solution}

\begin{solution}{iq:iv3:operating-systems-concurrency-and-networks:23}
\Cref{lst:iv3:os:framing} keeps a buffer, and takes out complete messages only. \texttt{recv} returns whatever has arrived, which may be part of a message, several messages, or a mixture. On loopback, a small write usually arrives whole, so the bad version passes; on a real network with segmentation and bursts, it does not.
\enlargethispage{-3\baselineskip}%
\omcode{code/interviews-3/05-operating-systems-concurrency-and-networks/python/iv_framing.py}{1}{26}{Framing a byte stream.\protect\hfil\protect\break}\label{lst:iv3:os:framing}

TCP has no idea where your messages begin. The length field goes first, so the reader can always tell how many bytes it needs, and the loop handles several messages in one read. If only a part of a message is there, it stays in the buffer until the next call.

The length limit matters. Without one, a corrupted length makes the reader wait for up to 65\,535 bytes that may never arrive, and then deliver garbage as a message, with every later frame misaligned. With a wider length field, the same fault could hold gigabytes. An empty body is a legal message, and the test covers it, as a careless \texttt{if not body} would lose one. The test also feeds 300 messages through random splits, and over a real TCP connection read in pieces of 97 bytes.

\texttt{del self.buf[:n]} is cheap in CPython, as a \texttt{bytearray} only moves its logical start, and the test sees the first byte's address move by $n$. In a language where it shifts the bytes, keep a read position, and compact from time to time.

On the sending side, \texttt{send} may accept only part of the data, so use \texttt{sendall} on a blocking socket, or keep the remainder and retry on a nonblocking one. Two threads must not write to one socket without a lock, or their messages can interleave.

\iqwrong{``TCP is reliable, so one \texttt{send} gives one \texttt{recv}.'' Reliability concerns the bytes. Nothing preserves the boundaries.}
\end{solution}

\begin{solution}{iq:iv3:operating-systems-concurrency-and-networks:24}
It prints \texttt{b\textquotesingle{}abcd\textquotesingle{} b\textquotesingle{}ij\textquotesingle{}}, then \texttt{65507}, then \texttt{Message too long}. The bytes \texttt{efgh} were discarded when the 4-byte read took the first datagram. Each \texttt{recv} consumes one whole datagram, whatever the buffer size. 65\,507 is the 65\,535 limit of the IP length field, minus 20 bytes of IP header and 8 of UDP.

Nothing tells you about the truncation by default. \texttt{recvmsg} sets a \texttt{MSG\_TRUNC} flag in its results, and with that flag passed to \texttt{recv}, the call returns the real length. So a feed handler should use a buffer larger than the biggest datagram it can receive, or check that flag.

A real feed keeps its datagrams below the path's MTU, typically under about 1500 bytes on Ethernet. Larger ones get split into IP fragments, and if one fragment is lost, the whole datagram is lost, so the loss rate multiplies. Fragments also add reassembly work. Venues specify their packet sizes, so read the specification and size the buffer to match.

On a loopback interface the MTU is large, which is why the maximum works here without any visible problem. That says nothing about a real network.

\iqwrong{``The rest of the datagram stays for the next \texttt{recv}.'' That is how a TCP stream behaves. In UDP, the next read gets the next datagram.}
\end{solution}

\begin{solution}{iq:iv3:operating-systems-concurrency-and-networks:25}
It prints \texttt{8192}, \texttt{sent 1000, no error}, \texttt{True} and \texttt{[True]}. The kernel doubles the requested size to allow for its own bookkeeping, so 4096 reads back as 8192. The receiving socket's buffer overflowed, the kernel dropped the rest there, and the last column of \texttt{/proc/net/udp} counts those drops.

The sender sees nothing, since UDP has no feedback. A successful \texttt{sendto} only means the packet was accepted for sending. Nothing tells the receiving program either, and that is why you must watch the counters. The count of datagrams received is below 1000 on any run where the buffer is this small and nobody reads while they are sent.

For a feed handler that drops during bursts, the order of action is as follows. First measure where: socket buffer drops, the NIC's ring, or the network. The sequence numbers tell you that data was lost, and the counters tell you where. Then look at the reader. It should do nothing on its thread but read and hand on, with slow work moved elsewhere. Give it a core of its own, and keep interrupts and the reader on the right NUMA node. Then raise \texttt{SO\_RCVBUF}, with the system limit \texttt{rmem\_max} if needed. A bigger buffer only rides out a burst, and if the reader is slower on average it still loses data, only later. For the last step, consider kernel bypass.

\iqwrong{``Make the buffer huge and the drops go away.'' It helps with bursts. It adds delay, and does nothing for a reader that is too slow overall.}
\end{solution}

\begin{solution}{iq:iv3:operating-systems-concurrency-and-networks:26}
A group is an address in 224.0.0.0/4 that any host can join. A sender, which need not join, sends one copy, and the network copies it to all members. A host joins with \texttt{IP\_ADD\_MEMBERSHIP}. With \texttt{join} the listing prints \texttt{b\textquotesingle{}tick\textquotesingle{}}, and otherwise \texttt{nothing}. 239.1.2.3 maps to \texttt{01:00:5e:01:02:03}, and 31 other groups share that address.
\omcode{code/interviews-3/05-operating-systems-concurrency-and-networks/python/iv_mcast_mac.py}{1}{7}{Group to Ethernet address.\protect\hfil\protect\break}\label{lst:iv3:os:mcastmac}

The Ethernet address keeps only the low 23 bits, and a group has 28 significant bits, so 5 bits are lost and $2^5 = 32$ groups map to each. 224.1.2.3 and 239.129.2.3 are two of them. So the network card may pass up packets for a group you did not join, and the kernel filters them out. Busy feeds should choose group addresses that do not collide.

Joining also has a network side. The host sends an IGMP report, and switches with IGMP snooping use it to forward the group only to ports that asked. The membership is per interface, so naming the wrong one is a classic reason for receiving nothing on a machine with several networks.

\texttt{IP\_MULTICAST\_TTL} defaults to 1, so the packet does not cross a router, and \texttt{IP\_MULTICAST\_LOOP} defaults to on, so a sender's own host receives it. We read both from a fresh socket. Therefore a test on one machine can work while two machines on different subnets get nothing, unless the TTL is raised and the routers are set up for multicast.

The run here is on loopback, with no switch and no IGMP on a wire, so what the network does rests on the sources in the ledger.

\iqwrong{``Binding to the group's port is enough.'' Without a join, the kernel has no reason to accept that group, as the second run shows.}
\end{solution}

\begin{solution}{iq:iv3:operating-systems-concurrency-and-networks:27}
It prints \texttt{(\textquotesingle{}EINPROGRESS\textquotesingle{}, 0)} and \texttt{(\textquotesingle{}EINPROGRESS\textquotesingle{}, \textquotesingle{}ECONNREFUSED\textquotesingle{})}. A nonblocking \texttt{connect} starts the handshake and returns at once, so its result is only ``in progress'', even when the attempt will fail. The real result is delivered later, in \texttt{SO\_ERROR}.

So writable does not mean connected. \texttt{select} marks the socket writable when the attempt has finished, in either way. If you take that for success and start sending, you see a confusing error on the first write, far from the cause. Read \texttt{SO\_ERROR} when it wakes: 0 means connected, and anything else is the reason for failure. Reading the option also clears it.

The listing's first case is a server that exists and the second is a closed port. Both return \texttt{EINPROGRESS} here, even on loopback, and that is what makes the check necessary. The manual allows other results, so handle a return of 0, which means it connected immediately, as well.

A blocking connect can wait a long time on an unreachable host, with several retransmissions of the SYN, so a process that must stay responsive connects without blocking, and sets its own time limit using \texttt{select} or \texttt{epoll}.

\iqwrong{``\texttt{select} said writable, so we are connected.'' It says the attempt has ended.}
\iqfollow{What should you do on failure?}{Close that socket and make a new one, as a failed connect leaves a socket in an unspecified state on some systems. Retry with a growing delay.}
\end{solution}

\begin{solution}{iq:iv3:operating-systems-concurrency-and-networks:28}
With no arguments, the C++ program is killed by \texttt{SIGPIPE} and prints nothing. With one, \texttt{send} uses \texttt{MSG\_NOSIGNAL}, so it returns $-1$ and the program prints \texttt{-1 Broken pipe}. In the Python listing, the first send returns 3, and the second raises \texttt{BrokenPipeError}, which is Python's name for \texttt{EPIPE}.

A process that writes to a connection with no reader gets a signal, and by default that signal ends the process. So a server can die because one client went away, unless it takes care. Set \texttt{SIGPIPE} to be ignored at start-up, or pass \texttt{MSG\_NOSIGNAL} on each send, and then handle the error return. Python sets it to ignored for you, which is why you see an exception there.

The first send succeeds, since at that point this side only knows that the peer has sent a FIN. A FIN means the peer has no more to send, and it could still be reading, so TCP accepts the data and sends it. The peer's kernel, finding that the socket has been closed and nobody can read this data, replies with RST. Once that has arrived, the next write fails. The listing waits for that reply with a \texttt{poll} that asks for no event, as only an error or a hang-up then wakes it. A \texttt{select} for reading would not wait, as the FIN has already made the socket readable. Across a network, with no wait, the second send could go first and also succeed.

So a successful \texttt{send} never meant that the peer read it. Bytes in the local buffer may be lost, and if you need to know, the application needs its own acknowledgement. For order entry, that is why a session has sequence numbers and replies.

\iqwrong{``A closed connection fails the next write at once.'' Only after the reset arrives, which takes a round trip.}
\end{solution}

\begin{solution}{iq:iv3:operating-systems-concurrency-and-networks:29}
The high-priority task finishes at 26. This is priority inversion: the task in the middle, which needs no lock, delays a higher one by all of its 20 ticks. With priority inheritance, the holder of the lock runs at its waiter's priority, and the high-priority task finishes at 6.
\omcode{code/interviews-3/05-operating-systems-concurrency-and-networks/python/iv_inversion.py}{1}{34}{One CPU and one lock.\protect\hfil\protect\break}\label{lst:iv3:os:inversion}

The trace without inheritance: the low task runs 0--1 and holds the lock. The high task arrives and blocks at once on the lock, so the low task carries on to 2, when the middle one arrives and takes the core until 22. Only then can the low task finish, at 24. The high task then takes 2 ticks. With inheritance, the low task runs through to 4 at high priority, and the high task finishes at 6, as if the middle task were not there.

The test shows the shape. Without inheritance, the high task finishes at $6 + m$ for a middle task of $m$ ticks, so the delay grows without limit. With it, the time is always 6, bounded by what is left of the critical section. This is checked on 200 random cases. The last completion time does not change with inheritance, as the total work is the same: it only changes who waits.

The other fixes: keep critical sections short, avoid sharing locks across priorities, or run all lock holders at a raised ceiling priority. Linux has \texttt{PTHREAD\_PRIO\_INHERIT} for mutexes. A trading system that uses real-time priorities on pinned cores has this same risk: a low-priority thread that holds a lock and has lost its core can stop the hot path. A \texttt{SCHED\_FIFO} thread that spins on a lock held by a lower real-time priority on its own core waits forever, as nothing lets the holder run. If the holder is a normal task, it runs only in the slice that the kernel keeps back from real-time tasks, 5\% by default.

\iqwrong{``Give the important task the highest priority, and it will always run first.'' Priority orders runnable tasks. A task that is blocked on a lock is not runnable.}
\end{solution}

\begin{solution}{iq:iv3:operating-systems-concurrency-and-networks:30}
\Cref{lst:iv3:os:rwlock} is one answer. \texttt{may\_read} is false once a writer is waiting. Without that preference, a steady flow of overlapping readers could hold off a writer forever. The risk simply moves: a steady flow of writers can now starve the readers.
\omcode{code/interviews-3/05-operating-systems-concurrency-and-networks/cpp/rwlock.hpp}{5}{42}{A reader-writer lock that prefers writers.\protect\hfil\protect\break}\label{lst:iv3:os:rwlock}

The last reader wakes a writer, as the writer's predicate cannot be true before that. A writer leaving wakes the next writer if there is one, and otherwise all readers, since many may enter at once. All state changes happen under the mutex, so no wake-up can be lost.

The test for preference needs no timing assumption in its assertion. One reader holds the lock. A writer thread calls \texttt{lock} and blocks, and the driver sleeps only to give it time to get there. Then \texttt{try\_lock\_shared} must return false. If the writer had not arrived yet, the call would return true and the test would fail, never pass wrongly. A twin with the writer check removed prints 1 on that line, so the test can fail the model. Under ThreadSanitizer, the writers and readers run with an invariant ($a = b$ on plain variables), and there is no race report.

It may be no faster than a mutex, since a reader must still write the lock's state, here taking an inner mutex. So on many cores, all readers contend for one cache line, and for a very short critical section that traffic is the main cost. A read lock pays off when readers hold it long enough for real overlap. We measured nothing here. Other options, for very short reads and rare writes, are a sequence lock, or replacing a pointer to an immutable copy.

\iqwrong{``Readers do not conflict, so a reader lock scales.'' Their work does not conflict. Their lock updates do.}
\end{solution}

\begin{solution}{iq:iv3:operating-systems-concurrency-and-networks:31}
\Cref{lst:iv3:os:futex} keeps 0 for free, 1 for locked, and 2 for locked with possible waiters. Locking and unlocking with no contention are atomic operations only, and the driver's million rounds make no \texttt{futex} call. Three states are needed so that \texttt{unlock} can skip the system call when nobody is asleep.
\omcode{code/interviews-3/05-operating-systems-concurrency-and-networks/cpp/futex_mutex.hpp}{7}{34}{A mutex on a futex.\protect\hfil\protect\break}\label{lst:iv3:os:futex}

With two states, \texttt{unlock} cannot know if anyone is sleeping, so it either makes a system call each time, or loses sleepers. State 2 records that someone may be asleep. A thread that fails to get the lock swaps in 2, and sleeps. If the old value was 0, its swap has in fact taken the lock.

A woken thread does not know whether others are still asleep, so it must set 2 and not 1. With 1, it would release later with no wake-up, and any other sleeper would sleep forever. Our test builds this twin, and the driver, with two waiters behind one holder, hangs and hits its time limit. The state 2 sometimes causes one unneeded wake-up, which is harmless.

\texttt{FUTEX\_WAIT} takes the expected value because of a race. Between a thread's decision to sleep and its entering the kernel, the holder may unlock and call wake, and with nobody yet asleep, that wake is lost. So the kernel checks the word and blocks atomically, and only if it still holds 2. If not, the call returns at once. That check closes the gap.

The memory orders are acquire on taking the lock, and release on giving it up, so the critical section's accesses cannot leak out. The counter's 400\,000 is exact, and ThreadSanitizer reports no race. A real mutex adds spinning before sleep, and a lot of careful handling, so use \texttt{std::mutex} in production.

\iqwrong{``A mutex always enters the kernel.'' Only when it has to sleep or to wake someone.}
\end{solution}

\begin{solution}{iq:iv3:operating-systems-concurrency-and-networks:32}
About 40.2 ms, or 400 times the round trip. The header goes at once, but Nagle holds back the body until the header is acknowledged, and the server delays that acknowledgement, as it has no reply to send until it has the body. Each waits for the other until the 40 ms timer ends it.
\omcode{code/interviews-3/05-operating-systems-concurrency-and-networks/python/iv_nagle.py}{1}{38}{Write, write, read, as events.\protect\hfil\protect\break}\label{lst:iv3:os:nagle}

Each rule is reasonable alone. Nagle's rule sends small data only when nothing is unacknowledged, which saved networks from floods of tiny packets. Delayed acknowledgement waits a little in the hope of carrying the ACK on a reply. The standard allows that delay up to 0.5 s, and the stem takes 40 ms. Together, in the pattern write, write, read, they deadlock until a timer fires. The timeline is: header arrives at 0.05, timer ends at 40.05, ACK arrives at 40.1, body arrives at 40.15, reply arrives at 40.2.

Three fixes. Best is to make one write: put the header and body in one buffer or use \texttt{writev}, which also saves a packet and a system call. Second is \texttt{TCP\_NODELAY}, which sends everything at once, for 0.1 ms in the model, and is the usual setting for trading connections, where the application decides what a message is. Third is \texttt{TCP\_QUICKACK} on the receiver, which gives 0.2 ms in the model, though it is not permanent, and must be set again.

One write alone never stalls, as the test shows, since nothing is outstanding when it goes. As the timer approaches 0, the stall shrinks to one extra round trip. \texttt{TCP\_NODELAY} defaults to 0 on a new socket, which we read back. This is a model, and no packet capture was run, so the real timer value is the machine's to determine.

\iqwrong{``Turning off Nagle makes TCP unreliable, or makes it stop retransmitting.'' It only stops holding back small segments.}
\end{solution}

\begin{solution}{iq:iv3:operating-systems-concurrency-and-networks:33}
The card takes the frame off the wire and writes it by DMA into a ring of receive buffers in memory. It raises an interrupt, or the driver is already polling. NAPI polls in software-interrupt context, hands the packet up through IP and UDP, finds the socket, and adds the packet to its receive queue. The scheduler wakes your thread, and \texttt{recv} copies the bytes to your buffer.

Stage by stage, here is where it waits, drops or copies.

The card: it hashes the header to choose one of its queues (receive side scaling), so a flow goes to a fixed queue and interrupt. If the ring is full, as when the kernel falls behind, the card drops the packet. \texttt{ethtool -S} shows those counters. Interrupt moderation may hold the interrupt for a short time, to batch packets.

The softirq: packets are processed on a core, possibly one chosen by the software steering of RPS. If that per-core queue passes its limit (\texttt{netdev\_max\_backlog}), the packet is dropped, and the second column of \texttt{/proc/net/softnet\_stat} on this kernel holds the count. If the core is busy with other work, the packet waits here.

The socket: UDP finds the socket by address and port, and multicast membership is checked. If the socket's receive buffer is full, the packet is dropped, and \texttt{/proc/net/udp} and \texttt{nstat} show that.

The wake-up: your thread must be scheduled, on a core that may be asleep or busy. A thread that spins in a nonblocking \texttt{recv} removes this wait. Then there's the copy to user space, one more copy, and the system call's cost.

Kernel bypass maps the card's rings directly into the process, and your thread polls them, so the interrupt, softirq, protocol stack, socket queue, wake-up and copy all go. What remains is the DMA, your poll, and a full ring if you are slow. You give up kernel tools, and have to run your own protocol handling.

This answer is a description, and no packet was traced for it, so the stages rest on the kernel documents in the ledger.

\iqwrong{``The packet goes from the card straight to my socket.'' There are several queues in between, and each one can drop.}
\end{solution}

\begin{solution}{iq:iv3:operating-systems-concurrency-and-networks:34}
It prints \texttt{[1, 1, 1]} and \texttt{[1, 0, 1]}. Level-triggered mode reports the socket as long as data is unread, so the six leftover bytes give an event on the second poll. Edge-triggered mode reports changes only: nothing new arrived, so the second poll is empty, and the third fires since new data came.

So under edge-triggered mode, if you read only part of the data and wait, you can wait forever with bytes in the buffer, if no more arrive. The rule is to read until \texttt{EAGAIN}. That needs a nonblocking socket, as a blocking read on an empty socket would stall the whole loop. The same goes for writes.

Edge-triggered mode avoids repeat wake-ups for a socket you know you cannot serve yet, and it lets several threads share a set of sockets without all waking up on one event. It is easier to get wrong: level-triggered is the safer default.

\texttt{select} and \texttt{poll} take the whole set of descriptors on each call, and the kernel scans it, so each call costs in proportion to the number watched, whatever the number active. \texttt{epoll} registers once and keeps a ready list, so each call costs in proportion to the number ready. For ten thousand idle connections, that is a large saving.

It is no real gain when nearly all descriptors are active on each call, or when there are only a handful of them. A hot path that watches a few sockets may not wait at all, but spin with nonblocking reads, and avoid the system call that \texttt{epoll\_wait} would be.

\iqwrong{``Edge-triggered is faster, so always use it.'' It avoids repeat events, but costs at least one extra read per burst, to see \texttt{EAGAIN}. Bugs in it hang silently.}
\end{solution}

\begin{solution}{iq:iv3:operating-systems-concurrency-and-networks:35}
It prints \texttt{parent threads 2}, \texttt{child threads 1, try\_lock 0} and \texttt{parent try\_lock 1}. Only the calling thread exists in the child. Memory is copied as it was, including the mutex in its locked state, and the thread that would unlock it does not exist there. So the child's lock is stuck forever.

\texttt{try\_lock} returns 0 here. A plain \texttt{lock} in the child would never return, and the program would hang with no error to show for it. That is the failure in the stem: a child calling into a logging library whose internal lock was held at the moment of the fork deadlocks at random, depending on what the other threads happened to be doing at the time. It may run for months and then hang.

The same goes for data under the lock. The other thread might have been halfway through an update, so the structure in the child is inconsistent, and no thread is there to finish it. Some libraries guard their state with \texttt{pthread\_atfork} handlers, but you cannot count on every library.

So until it calls \texttt{exec}, a child of a multi-threaded process may safely call only async-signal-safe functions. In practice, you do almost nothing and call \texttt{exec} at once, or use \texttt{posix\_spawn}. If you fork for a snapshot, do it before starting any threads, or copy data out by another route.

\iqwrong{``The child gets all the threads, as they are part of the process.'' It gets one. Everything the others were doing just stops, in the child's copy.}
\end{solution}

\begin{solution}{iq:iv3:operating-systems-concurrency-and-networks:36}
With no arguments it prints \texttt{EOWNERDEAD} and \texttt{0}: the lock is granted, along with the news that its owner died. With one argument, a plain mutex stays locked, and the timed lock gives \texttt{ETIMEDOUT}. Without a time limit, the survivor would hang forever, with nothing in the logs.

The kernel can do this since a robust mutex's owner keeps a list of the robust locks it holds. When a thread dies, the kernel walks that list, marks each lock with a flag, and wakes a waiter. A plain mutex is only a word in memory, and nobody knows who holds it. So a process crash leaves that word locked, and the whole shared segment is stuck.

On \texttt{EOWNERDEAD}, the new owner holds the lock, and the data it guards may be half updated, as the owner died at an arbitrary instruction. Only the application knows what a correct state looks like. So the new owner must check or repair the data, and then call \texttt{pthread\_mutex\_consistent}. If it unlocks without that call, the mutex becomes permanently unusable, and all later lock calls return \texttt{ENOTRECOVERABLE}. The library cannot repair the data for you, since it does not know what the lock protects.

So robustness has a price: every update must leave the data repairable, for instance by writing a log first, or by marking a record dirty during the update. In practice, the data's design is the hard part.

The other course is to avoid a lock across processes at all. A lock-free ring lets a producer die with no lock held, and a separate process can watch for the death with \texttt{waitpid} or a heartbeat.

\iqwrong{``The kernel frees a dead process's locks.'' It frees its memory and files. A plain mutex in shared memory is nobody's resource.}
\end{solution}

\begin{solution}{iq:iv3:operating-systems-concurrency-and-networks:37}
\Cref{lst:iv3:os:arb} keeps the next number expected and a set of later ones. A duplicate or an old number is dropped, and whatever becomes contiguous comes out. When more than \texttt{hold} are waiting, the missing range is reported and skipped. Waiting forever would freeze the book: if both lines lost a packet, it will never come.
\omcode{code/interviews-3/05-operating-systems-concurrency-and-networks/python/iv_arb.py}{1}{26}{Merging two lines.\protect\hfil\protect\break}\label{lst:iv3:os:arb}

The first copy wins, so the merged stream is at least as fast as the faster line, packet by packet, which is the point of having two. Each line is usually in order, so a missing number mostly means that line lost it, and the other line's copy arrives shortly after.

\texttt{hold} sets the trade-off. Too small, and a copy that was merely behind on the slower line gets declared lost, which triggers a recovery for nothing. Too large, and every real loss stalls all later data until it is recognised. Choose it from the measured difference between the two lines: how many packets can arrive on the faster line during the normal lag of the slower one, at peak rates. A real handler adds a timer, as during a quiet market no later packet comes to trigger the gap check. This version has no timer, so a gap at the end of a burst stays pending.

The declared gap must go somewhere: a retransmission request, or a snapshot from a recovery service. Late copies that arrive after the skip are then dropped, as they are older than \texttt{next}.

The test merges two independently lossy lines at random on 300 cases. The output must be increasing, with no duplicates, and all numbers lost on both lines and passed by the output must be reported as gaps. With a large \texttt{hold}, a loss on one line never becomes a gap.

\iqwrong{``Forward whatever arrives, and let the book sort it out.'' A book applied out of order gives wrong prices, so a single stream in sequence must come out of this layer.}
\end{solution}

\begin{solution}{iq:iv3:operating-systems-concurrency-and-networks:38}
Each connection has a read buffer, through its \texttt{Deframer}, and a write buffer, \texttt{out}. Pending replies go in \texttt{out}. \cref{lst:iv3:os:loop1,lst:iv3:os:loop2} ask for \texttt{EPOLLOUT} only while \texttt{out} holds data, as an idle socket is almost always writable, and the loop would wake without pause. The server never blocks on one client.
\omcode{code/interviews-3/05-operating-systems-concurrency-and-networks/python/iv_event_loop.py}{7}{38}{The event loop.\protect\hfil\protect\break}\label{lst:iv3:os:loop1}
\enlargethispage{-2\baselineskip}%
\omcode{code/interviews-3/05-operating-systems-concurrency-and-networks/python/iv_event_loop.py}{41}{68}{Reading and sending for one connection.\protect\hfil\protect\break}\label{lst:iv3:os:loop2}

The loop uses the default level-triggered mode, so one \texttt{recv} per event is enough, and leftover data gives another event. Every socket is nonblocking, and that matters most for \texttt{send}: on a blocking socket, a client that stopped reading would stop the server for everyone, once its buffer filled.

\texttt{send} may take only part of the data, so \texttt{flush} removes what was accepted and keeps the rest. When \texttt{recv} returns empty, the peer has finished sending, and it may still be reading, as with a half-close. So the server stops reading, sends what is in \texttt{out}, and then closes. If the peer has really gone, a \texttt{recv} or a \texttt{send} raises an error, and the connection is dropped with its pending replies, which have nowhere to go. A bad length from one client closes that client alone.

In the test, 8 clients each send 300 messages of up to 4000 bytes, chopped at random, then half-close, and they read in pieces of 1000 bytes on a separate thread. The sockets have small buffers, so the server must queue replies, and the test checks that it did. Every client gets its messages back exactly, in order, and the server stops once all have closed. A second test has one client close without reading and another send a length above the limit, and a third client is still served.

One slow client costs the others little CPU, but this listing has a real flaw: \texttt{out} has no limit, so a client that writes and never reads makes the server's memory grow without bound. A real server limits it, by no longer reading from that client once \texttt{out} is large, which pushes back through TCP to the sender, or by closing the connection. Also, a single huge message delays the rest, as the loop runs one thread. Handlers must not block or run long.

\iqwrong{``Register \texttt{EPOLLOUT} for all sockets from the start.'' With level-triggered mode, the loop spins at full CPU.}
\end{solution}

\begin{solution}{iq:iv3:operating-systems-concurrency-and-networks:39}
\Cref{lst:iv3:os:reliable} is one such model. The sender repeats every unacknowledged packet in its window, and the receiver holds early packets, and delivers contiguous ones. It must acknowledge a duplicate, as a duplicate means its earlier ACK may have been lost. Otherwise the sender would resend forever.
\omcode{code/interviews-3/05-operating-systems-concurrency-and-networks/python/iv_reliable.py}{6}{43}{Selective repeat, in rounds.\protect\hfil\protect\break}\label{lst:iv3:os:reliable}

The window limits how much is in flight, so it limits how much the receiver may have to hold and how much a loss wastes. With $w = 1$, this is stop-and-wait, and the rate is one packet per round trip, which the test checks through the widest window used.

Why in order and exactly once? \texttt{delivered} grows only in the \texttt{while} loop, which appends \texttt{next} and increments it, so the numbers delivered are $0, 1, 2, \dots$ with none skipped or repeated. A packet with a number below \texttt{next} is never stored, so a duplicate cannot come out again. Why does it finish? While the loss rate is below 1, each packet eventually gets through with its ACK, with probability 1, so the base keeps moving. The sender's \texttt{acked} set only ever holds packets that the receiver really received. So nothing is dropped by a false acknowledgement.

The test delivers all of 200 packets at loss rates of 0, 0.1, 0.5 and 0.8, on several seeds. It checks the window, that no loss means exactly 200 transmissions, that more loss costs more, and that a loss rate of 1 raises \texttt{TimeoutError} rather than hanging. Duplicates and reordering are in the link model.

TCP has the same core, a window with cumulative acknowledgements, and selective ones as an option. What the model leaves out is real time: retransmission timers, round-trip estimation, and a congestion window that shrinks on loss. Flow control is also absent: this receiver has unlimited buffer. In a real network, sending a full window on each round when the link is the cause of the loss makes it worse.

\iqwrong{``Number the packets, and retransmit when the receiver asks.'' If the request itself is lost, nothing happens. Some side needs a timer.}
\end{solution}

\begin{solution}{iq:iv3:operating-systems-concurrency-and-networks:40}
\Cref{lst:iv3:os:shmring} is the ring. Indices must be lock-free, as a lock inside \texttt{std::atomic} would be private to one process, or be stuck if its holder died. The \texttt{static\_assert} on \texttt{is\_always\_lock\_free} checks that at compile time. If the producer dies, the consumer sees everything that was published, each slot whole, and never a slot half written.
\omcode{code/interviews-3/05-operating-systems-concurrency-and-networks/cpp/shm_ring.hpp}{6}{27}{A ring in shared memory.\protect\hfil\protect\break}\label{lst:iv3:os:shmring}

The slot is written first, and \texttt{head} moves only afterwards, with a release store. So a kill at any instruction leaves one of two states: \texttt{head} unchanged, and the slot's contents simply unseen, or \texttt{head} moved with the slot complete. One store to \texttt{head} is one atomic step, with no state between. The consumer cannot tell from the ring alone that the producer has died, as a dead one and an idle one look the same, so it needs \texttt{waitpid}, a heartbeat or similar.

The driver has a child push the sequence 0, 1, 2 and so on until the parent kills it with \texttt{SIGKILL}, at a moment it did not choose. The parent has checked 200\,000 values by then, drains what remains, and finds the sequence unbroken. It prints \texttt{0 9 1}: no errors, signal 9, and enough read.

A mutex in that memory would behave differently. A process killed while holding one leaves it locked, unless it is robust, and then the data can be half updated, so you need a recovery plan. For this reason, a shared-memory path between processes is often built lock-free.

ThreadSanitizer sees one process, so the test runs the same header between two threads under it, and finds no race. With the store and load relaxed, it reports a \texttt{data race}, which shows that the test can catch a wrong order. The orders are the same, whether threads or processes run them. x86-64 hides this kind of fault on real runs, and we had no weakly ordered machine here. The cache-line alignment keeps the two indices from sharing a line.

\iqwrong{``After a crash, the consumer may read a half-written message.'' Only if \texttt{head} moves first. The order of the two writes is the whole design.}
\end{solution}
